\section{Covariance and physical relative Hessian vertices}
\label{chr-section}

This section controls the first two background coefficients of an
actual action Hessian.  It does not estimate the moving hard
mean-zero projection, nor compute a logarithmic determinant
coefficient.  The argument uses exact gauge covariance before
taking absolute values.

\subsection{Hypotheses and conclusion}
\label{chr-hypotheses}

Work on the current unit lattice with counting inner products and a
fixed finite bond/color fiber.  Let \(L(a)\) be a transpose-symmetric
family of left-relative bond Hessians.  Assume that it is exactly
covariant under sitewise adjoint frames and that its current-coordinate
Taylor kernels through background degree two at zero have uniform
exponentially rooted absolute row bounds.  The transposed bounds are
included.  For an action Hessian, these are the action kernels through
degree four together with the local relative-coordinate changes.

Write \(q=L(0)\).  Assume that \(q\) is translation invariant, scalar
on each color factor, and exponentially summable, uniformly in the
old classical depth.  In source-site coordinates impose the
componentwise moment conditions
\begin{equation}
 \sum_r q_{\alpha\beta}(r)=0,\qquad
 \sum_r r_\rho q_{\alpha\beta}(r)=0
 \quad\text{for every }\alpha,\beta,\rho .
 \label{chr-moments}
\end{equation}
Thus its Fourier symbol vanishes to order two.  No symmetry of
\(q_{\alpha\beta}(r)\) in \(\alpha,\beta\) at a fixed \(r\) is assumed.
Only the full transpose relation
\(q_{\alpha\beta}(r)=q_{\beta\alpha}(-r)^T\) is used.

The actual application here is on a chosen infinite structural system,
where these continuous-momentum moments have been established.
Explicit periodizations of that system are also covered with the
specified lifted paths.  An arbitrary actual finite-regulator winding
kernel is not silently identified with such a periodization.

Set \(\mu=K^{-1}\leq1\).  For a complete physical field \(A\), define
\begin{equation}
 \|A\|_{\mathrm{phys},1}
 =\|A\|_\infty+
   \mu^{-1}\max_{\substack{e,e'\text{ nearest}\\
                          \text{same component}}}
                     |A(e')-A(e)| .
 \label{chr-physical-norm}
\end{equation}
Let \(L_j=[t^j]L(t\mu A)\), \(j=1,2\), and suppose
\(\|A\|_{\mathrm{phys},1}\leq a_*\).
Here \(D_0\) denotes the ordinary componentwise gradient of a
\emph{bond field}, not the scalar gauge generator.  Put
\(R_0=D_0^TD_0+\mu^2I\).

\begin{theorem}
\label{chr-main}
Under these hypotheses there are operators \(Z_j,E_j\) such that
\begin{equation}
 L_j=D_0^TZ_j+Z_j^TD_0+E_j,\qquad j=1,2,
 \label{chr-decomposition}
\end{equation}
with
\begin{align}
 \|Z_1\|&\leq C\mu a_*,
 &\|E_1\|&\leq C\mu^2a_*,
 \notag\\
 \|Z_2\|&\leq C\mu^2a_*^2,
 &\|E_2\|&\leq C\mu^2a_*^2 .
 \label{chr-bounds}
\end{align}
Consequently
\begin{equation}
 \|R_0^{-1/2}L_jR_0^{-1/2}\|
       \leq C_j a_*^j,\qquad j=1,2.
 \label{chr-relative}
\end{equation}
The constants are independent of \(K\) and of the old depth.
Mixed second coefficients follow by polarization.  For directions
supported in fixed smooth physical collars, the factors admit
exponentially rooted final-block operator bounds with the same
powers of \(\mu\).
\end{theorem}

For the actual homogeneous classical action \(S_M\), the hypotheses
come from the established uniform support and Cauchy bounds, gauge
covariance, and the flat normalization
\[
 q_{M,\mathrm{mid}}(p)=|p|^2I-pp^T+O(|p|^4).
\]
The analytic midpoint-to-source-site phases preserve the zero and
first jets, so they imply \eqref{chr-moments}.  This last input is
stronger than row-sum cancellation alone.

\subsection{Formal local path gauges}
\label{chr-local-gauge}

For each displacement \(r\), choose a translation-covariant
coordinate path of length \(|r|_1\), with the negative-displacement
path defined by reversal.  Denote its adjoint transport by
\(W_{xy}(a)\).  It obeys \(W_{yx}=W_{xy}^T\) and is complex
orthogonal coefficientwise.

Fix the source site \(x\) of a matrix row.  The transports from \(x\)
to every \(z\) define a formal site gauge \(g_x(z)\), with
\(g_x(x)=I\).  A transformed edge is a loop consisting of the two
selected paths and that edge.  Its length is at most
\(C(1+|z-x|)\); a plaquette filling uses at most
\(C(1+|z-x|)^2\) plaquettes.  Write its link-angle jet as
\[
 b(t,z)=t b_1(z)+t^2b_2(z)+O(t^3).
\]
Discrete Stokes at first order and the physical first-difference
bound give
\begin{equation}
 |b_1(z)|\leq C\mu^2a_*(1+|z-x|)^2.
 \label{chr-b1}
\end{equation}
The second coefficient of the product and its logarithm is a sum
of quadratic pairs of first link coefficients.  There are at most
the square of the loop length such pairs.  Hence
\begin{equation}
 |b_2(z)|\leq C\mu^2a_*^2(1+|z-x|)^2.
 \label{chr-b2}
\end{equation}

These are formal coefficient identities.  The potentially
unbounded gauge \(g_x\) is not placed in a global analytic ball.
One first truncates its polynomial jets far from \(x\), applies
the finite covariance identity, and then passes coefficientwise.
The exponentially rooted bounds absorb the polynomial tails.
In particular this step does not assert a finite-period identity
for an unbounded affine field.

Gauge covariance and \(g_x(x)=I\) give, with a consistent choice
of transport orientation,
\begin{equation}
 L(t\mu A)_{x\alpha,y\beta}
        =L(b(t))_{x\alpha,y\beta}W_{xy}(t\mu A).
 \label{chr-gauge-identity}
\end{equation}
The inverse convention gives the same argument with both
orientations reversed.  Define the dressed flat kernel
\[
 L^P(t)_{x\alpha,y\beta}
     =q_{\alpha\beta}(y-x)W_{xy}(t\mu A).
\]
Using \eqref{chr-b1}--\eqref{chr-b2} and the rooted Taylor kernels,
\begin{align}
 \|[t](L-L^P)\|_{\mathrm{Schur}}
     &\leq C\mu^2a_*,
 \notag\\
 \|[t^2](L-L^P)\|_{\mathrm{Schur}}
     &\leq C\mu^2a_*^2 .
 \label{chr-peierls-error}
\end{align}
At second order the terms are \(L'[b_2]\),
\(\frac12L''[b_1,b_1]\), and the product of \(L'[b_1]\)
with the first transport coefficient.  The latter two have,
respectively, the stronger powers \(\mu^4\) and \(\mu^3\).
All sums are fine current-coordinate Schur sums; no final-block
fiber dimension is inserted.

\subsection{Ordered factorization and covariant completion}
\label{chr-factorization}

A common difference-square formula would be invalid here: the
product of two equal-displacement differences is symmetric in
the bond components and would discard the odd-displacement,
skew-component part of \(q(r)\).  Even the forward-gradient Maxwell
kernel has such shifted entries.  Instead use an ordered Laurent
factorization.

For a Laurent series \(F\), define
\[
 {\cal B}_jF(z)=
 \frac{F(1,\ldots,1,z_j,\ldots,z_d)
       -F(1,\ldots,1,1,z_{j+1},\ldots,z_d)}{z_j-1}.
\]
The quotient is removable.  Applying the telescoping identity
twice, and using \eqref{chr-moments}, gives
\begin{equation}
 q(z)=\sum_{\rho,\tau}(z_\rho^{-1}-1)
       H_{\rho\tau}(z)(z_\tau-1),\qquad
 H_{\rho\tau}(z)=-z_\rho{\cal B}_\tau{\cal B}_\rho q(z).
 \label{chr-laurent}
\end{equation}
Each monomial division costs at most \(|r|_1\) terms with no
increase in exponent length.  Spending an exponential reserve
twice therefore yields a uniform exponentially summable \(H\).
For example an input exponent \(\sigma\) gives an output
exponent \(\sigma/2\) and a bound
\[
 \|H\|_{\sigma/2}
 \leq \frac{16e^{\sigma/2}}{e^2\sigma^2}\|q\|_\sigma
\]
in the combined absolute row norm; the analogous column bound
follows in the same way.  Replace \(H\) by \((H+H^T)/2\), where
the transpose includes displacement reversal and both gradient
and bond indices.  The factorization is unchanged.  We have
\begin{equation}
                  q=D_0^TH_0D_0,\qquad H_0=H_0^T.
 \label{chr-flat-factor}
\end{equation}
The color-scalar property is preserved.

Dress each \(H_0\) entry by its endpoint transport:
\[
 H^P(a)_{x\rho\alpha,y\tau\beta}
       =(H_0)_{\rho\alpha,\tau\beta}(y-x)W_{xy}(a).
\]
Let \(D_a\) be the nearest-neighbor covariant component gradient
and define
\begin{equation}
                  Q_{\mathrm{cov}}(a)
                   =D_a^TH^P(a)D_a .
 \label{chr-completion}
\end{equation}
This family is exactly covariant and transpose symmetric, with
\(Q_{\mathrm{cov}}(0)=q\).  Its first two Taylor kernels have
uniform exponentially rooted bounds: path derivatives only add
polynomial length factors.  On a sufficiently small complex
current-angle ball, the path-product factor
\(\exp(C\|a\|_\infty|r|)\) is paid by the exponential reserve.

Apply the local-gauge argument to \(Q_{\mathrm{cov}}\) as well
as to \(L\).  They have the same flat \(q\), hence the same dressed
flat comparison \(L^P\).  Subtraction of their two estimates
\eqref{chr-peierls-error} gives symmetric errors
\begin{equation}
 {\cal E}_j=[t^j](L(t\mu A)-Q_{\mathrm{cov}}(t\mu A)),
 \qquad
 \|{\cal E}_j\|_{\mathrm{Schur}}
           \leq C\mu^2a_*^j,\quad j=1,2.
 \label{chr-completion-error}
\end{equation}
This argument does not presume that the paths generated by
\eqref{chr-completion} coincide with the initially selected
endpoint paths.

\subsection{Coefficient expansion and localization}
\label{chr-expansion}

Write, through second order,
\[
 D_{t\mu A}=D_0+tF_1+t^2F_2,\qquad
 H^P(t\mu A)=H_0+tH_1+t^2H_2.
\]
Nearest-neighbor differentiation and exponential path moments give
\begin{align}
 \|F_1\|+\|H_1\|&\leq C\mu a_*,
 \notag\\
 \|F_2\|+\|H_2\|&\leq C\mu^2a_*^2 .
 \label{chr-factor-bounds}
\end{align}
All \(H_j\) are symmetric.  Set
\begin{align}
 Z_1&=\tfrac12H_1D_0+H_0F_1,
 \notag\\
 Z_2&=\tfrac12H_2D_0+H_1F_1+H_0F_2,
 \notag\\
 E_1&={\cal E}_1,\qquad
 E_2=F_1^TH_0F_1+{\cal E}_2 .
 \label{chr-exact-z}
\end{align}
Direct expansion of \eqref{chr-completion} proves
\eqref{chr-decomposition}.  The counting norm of \(D_0\) is bounded
by a dimensional constant, so \eqref{chr-factor-bounds} and
\eqref{chr-completion-error} prove \eqref{chr-bounds}.
Finally,
\[
 \|D_0R_0^{-1/2}\|\leq1,\qquad
 \|R_0^{-1/2}\|\leq\mu^{-1},
\]
which proves \eqref{chr-relative}.

For localization, use complete directions supported in a fixed
smooth physical partition with bounded overlap and bounded physical
first derivatives.  The terms \(F_1,F_2\) have nearest-neighbor
support.  Each marked source in \(H_1,H_2\) lies on a selected
endpoint path.  A distant source forces a long path, and the
exponential \(H_0\) kernel pays its length.  In the remainders,
a source on a local-gauge path is connected to the fixed matrix
row through the original background-source endpoint.  The
original rooted Taylor bound pays that distance.  Discrete
Stokes uses source differences in the same smooth collar,
enlarged by one fine edge.  Polynomial path multiplicities are
absorbed by a smaller exponential exponent.

Thus the factors in \eqref{chr-exact-z} and the \(E_j\) have
rooted final-block bounds of the form
\(\exp[-c\tau(\text{row},\text{column},\text{source labels})]\)
with the powers in \eqref{chr-bounds}.  Fine-coordinate Schur
summation proves the block-operator estimate; summing \(K^4\)
independent entry bounds is unnecessary.  No sharp source cutoff
is differentiated.

Physical weights \(\exp(\theta f(x/K))\), with \(f\) Lipschitz,
change a path factor by at most
\(\exp(C\theta\mu|r|)\), again paid by the exponential reserve.
Commuting \(D_0\) past the weight produces \(O(\theta\mu)\).
When multiplied by \(Z_1\) or \(Z_2\), this has the zero-order
size required for \(E_1\) or \(E_2\).  This supplies the weighted
localized form of the theorem.

\subsection{Scope of the actual application}
\label{chr-scope}

For \(L=S_M''\), the hypotheses are the uniform action support and
Cauchy bounds through degree four, actual gauge covariance, and
the already normalized flat kernel.  The result concerns bounded
physical \(C^1\) fields and their localized jets.  It does not give
the same gain for rapidly oscillating bounded potentials, whose
fine curvature may be \(O(K^{-1})\), and it gives no estimate for
the hard-projection scalar operator.

In the intended soft-vector application, the exact on-shell
identification and its source-coordinate conventions must first
justify the relevant affine jets of
\[
 S_M''(e^{iA/K})+D_AD_A^T+
 \overline C_A^T
 \bigl[aI-S_{MK}''(M_K(e^{iA/K}))\bigr]\overline C_A .
\]
The corrected constraint's full \(O(K^{-1})\) bound and its
physical source derivatives are separate inputs.  The present
theorem supplies the action-Hessian relative vertices.  It does
not by itself identify a soft-vector determinant, control every
stationary response, or determine a one-loop sign.

